\chapter{参数估计}

\section{点估计}

点估计：用统计量的一个具体数值 $\hat{\theta}$ 来估计未知参数 $\theta$。

常用的构造估计量的方法有矩估计法和极大似然估计法。

\subsection{矩估计法}

\textbf{核心思想：} 用样本矩代替（估计）总体矩。

设总体 $X$ 的分布含 $k$ 个未知参数 $\theta_1, \dots, \theta_k$：
\begin{align}
    \mu_\ell = \E[X^\ell] &= g_\ell(\theta_1, \dots, \theta_k), \quad \ell = 1,\dots,k.
\end{align}

\textbf{步骤：}
\begin{enumerate}
    \item 计算总体矩 $\mu_\ell = \E[X^\ell]$（用参数表示）；
    \item 令样本矩等于总体矩：$A_\ell = \mu_\ell$；
    \item 解方程组得到 $\hat{\theta}_1, \dots, \hat{\theta}_k$。
\end{enumerate}

\begin{example}[正态分布的矩估计]
设 $X \sim \mathcal{N}(\mu, \sigma^2)$：
\begin{align}
    \E[X]    &= \mu, \\
    \E[X^2]  &= \Var(X) + (\E[X])^2 = \sigma^2 + \mu^2.
\end{align}
令 $A_1 = \mu$，$A_2 = \sigma^2 + \mu^2$，解得：
\begin{equation}
    \hat{\mu} = \bar{X},\quad \hat{\sigma}^2 = A_2 - A_1^2 = \frac{1}{n}\sum_{i=1}^{n}(X_i - \bar{X})^2.
\end{equation}
\end{example}

\subsection{极大似然估计 (MLE)}

\begin{definition}[似然函数]
设 $x_1, \dots, x_n$ 为来自总体 $X$ 的样本观测值，总体的分布为 $f(x; \theta)$（离散时为概率质量函数，连续时为概率密度函数），则\textbf{似然函数}为：
\begin{equation}
    L(\theta) = L(\theta; x_1, \dots, x_n) = \prod_{i=1}^{n} f(x_i; \theta).
\end{equation}
\end{definition}

\begin{remark}[符号 $L(\theta; x_1,\dots,x_n)$ 中分号的含义]
$L(\theta; x_1, \dots, x_n)$ 中的分号 \textbf{不是}条件概率的竖线 $|$，它只是一个分隔符，
用于区分两类量：

\begin{itemize}
    \item \textbf{分号前}（$\theta$）：当前关注的\textbf{变量}——MLE 中我们要对它求导、最大化
    \item \textbf{分号后}（$x_1,\dots,x_n$）：已观测到的\textbf{固定数据}——拿到手后就不再改变
\end{itemize}

写成 $L(\theta \mid x_1,\dots,x_n)$ 也不少见，但分号写法在统计学中更规范，因为：

\begin{center}
\renewcommand{\arraystretch}{1.3}
\begin{tabular}{c|c|c}
\toprule
\textbf{记号} & \textbf{含义} & \textbf{例子} \\
\midrule
$f(x; \theta)$ 或 $p(x \mid \theta)$ & $\theta$ 固定时，$x$ 的概率密度/质量函数 & 正态 PDF \\
$L(\theta; x)$ 或 $L(\theta \mid x)$ & $x$ 固定时，$\theta$ 的似然函数 & MLE 的目标函数 \\
$P(\theta \mid x)$ & 观测到 $x$ 后，$\theta$ 的\textbf{后验概率}（贝叶斯） & 不是 MLE！需要先验 \\
\bottomrule
\end{tabular}
\end{center}

\textbf{关键区分：} $L(\theta; x)$ 和 $P(\theta \mid x)$ 是完全不同的量。
前者是频率学派的核心工具（MLE），后者是贝叶斯学派的核心工具（MAP / 后验推断）。
两者通过贝叶斯公式关联：$P(\theta \mid x) \propto L(\theta; x) \cdot P(\theta)$。
\end{remark}

\textbf{核心思想：} 选择使观测数据出现概率最大的参数值。

极大似然估计值 $\hat{\theta}$ 满足：
\begin{equation}
    L(\hat{\theta}) = \max_{\theta \in \Theta} L(\theta).
\end{equation}

由于 $\ln$ 单调递增，通常最大化\textbf{对数似然函数}：
\begin{equation}
    \ell(\theta) = \ln L(\theta) = \sum_{i=1}^{n} \ln f(x_i; \theta).
\end{equation}

求 $\dfrac{\partial \ell}{\partial \theta} = 0$ 解得 $\hat{\theta}_{\mathrm{MLE}}$。

\begin{remark}[为什么 MLE 常写成 $\argmax_\theta P(\mathcal{D} \mid \theta)$？这不是贝叶斯]
MLE 的两种等价写法常引起混淆：
\[
\hat{\theta}_{\mathrm{MLE}} = \argmax_\theta L(\theta; \mathcal{D})
                           = \argmax_\theta P(\mathcal{D} \mid \theta).
\]

\textbf{$P(\mathcal{D} \mid \theta)$ 不是贝叶斯，因为它读作"给定参数，数据的概率"。}

竖线 $|$ 在概率论中是\textbf{条件概率}的标准符号：$P(A \mid B) = P(A \cap B) / P(B)$。
$P(\mathcal{D} \mid \theta)$ 的含义是："当参数取值为 $\theta$ 时，数据 $\mathcal{D}$ 的分布"——这正是采样分布/似然。
频率学派和贝叶斯学派\emph{都使用} $P(\mathcal{D} \mid \theta)$，分歧在于用它做什么：

\begin{center}
\renewcommand{\arraystretch}{1.4}
\begin{tabular}{c|c|c}
\toprule
& \textbf{频率学派 (MLE)} & \textbf{贝叶斯学派} \\
\midrule
用的量          & $P(\mathcal{D} \mid \theta)$ & $P(\mathcal{D} \mid \theta) \cdot P(\theta)$ \\
对 $\theta$ 的立场 & 固定未知常数，放在竖线右边作为"条件" & 随机变量，有先验分布 $P(\theta)$ \\
目标            & 找使 $P(\mathcal{D} \mid \theta)$ 最大的 $\theta$ & 求 $P(\theta \mid \mathcal{D}) \propto P(\mathcal{D} \mid \theta)P(\theta)$ \\
\bottomrule
\end{tabular}
\end{center}

\textbf{分号 $;$ 和竖线 $|$ 在这里可以互换吗？}

\begin{itemize}
    \item $L(\theta; x)$ 用分号强调"视角翻转"：数据固定，$\theta$ 是变量
    \item $P(x \mid \theta)$ 用竖线因为它确实满足条件概率的数学定义（$\theta$ 固定时 $x$ 的分布）
    \item 两者数学上完全相等：$L(\theta; x) = P(x \mid \theta) = \prod f(x_i; \theta)$
    \item 分号更"安全"（不会误以为是贝叶斯），竖线更"标准"（符合概率论规范）。两种写法都正确
\end{itemize}

\textbf{MLE 从未把 $\theta$ 当作随机变量。} $P(\mathcal{D} \mid \theta = 0.5)$ 的意思是"假设硬币正面概率恰好是 0.5，那么掷出 8 正 2 反的概率是多少"——"$\theta = 0.5$"只是一个假定的场景，不是随机事件。这就像"假如明天下雨，你带伞的概率是多少"——"下雨"在条件位置，不意味着你把天气当作随机变量来建模。

只有当竖线左右颠倒——$P(\theta \mid \mathcal{D})$ ——才是贝叶斯，因为此时参数被当作随机变量。
\end{remark}

\subsubsection{完整示例：正态分布 $N(\mu, \sigma^2)$ 的 MLE}

设 $X \sim \mathcal{N}(\mu, \sigma^2)$，$\mu, \sigma^2$ 均未知。$x_1, x_2, \dots, x_n$ 为样本观测值。求 $\mu, \sigma^2$ 的 MLE。

\medskip
\textbf{第 1 步：写出 PDF，构造似然函数。}

正态分布的 PDF：
\begin{equation}
    f(x; \mu, \sigma^2) = \frac{1}{\sqrt{2\pi\sigma^2}} \exp\!\left(-\frac{(x - \mu)^2}{2\sigma^2}\right).
\end{equation}

将 PDF 中的参数固定位置换成变量视角，似然函数为 $n$ 个独立观测的联合密度：
\begin{equation}
    L(\mu, \sigma^2) = \prod_{i=1}^{n} f(x_i; \mu, \sigma^2)
    = \prod_{i=1}^{n} \frac{1}{\sqrt{2\pi\sigma^2}} \exp\!\left(-\frac{(x_i - \mu)^2}{2\sigma^2}\right).
\end{equation}

\medskip
\textbf{第 2 步：取对数，化积为和。}

\begin{align}
    \ell(\mu, \sigma^2) = \ln L(\mu, \sigma^2)
    &= \sum_{i=1}^{n} \left[ -\frac{1}{2}\ln(2\pi) - \frac{1}{2}\ln(\sigma^2) - \frac{(x_i - \mu)^2}{2\sigma^2} \right] \\
    &= -\frac{n}{2}\ln(2\pi) - \frac{n}{2}\ln(\sigma^2) - \frac{1}{2\sigma^2}\sum_{i=1}^{n}(x_i - \mu)^2.
\end{align}

\medskip
\textbf{第 3 步：分别对 $\mu$ 和 $\sigma^2$ 求偏导，令其为零。}

对 $\mu$ 求偏导：
\begin{equation}
    \frac{\partial \ell}{\partial \mu}
    = -\frac{1}{2\sigma^2} \sum_{i=1}^{n} 2(x_i - \mu)(-1)
    = \frac{1}{\sigma^2} \sum_{i=1}^{n} (x_i - \mu) = 0.
\end{equation}

解得：
\begin{equation}
    \boxed{\hat{\mu} = \frac{1}{n}\sum_{i=1}^{n} x_i = \bar{x}}.
\end{equation}

对 $\sigma^2$ 求偏导（注意是对 $\sigma^2$ 整体，而非 $\sigma$）：
\begin{equation}
    \frac{\partial \ell}{\partial \sigma^2}
    = -\frac{n}{2} \cdot \frac{1}{\sigma^2} + \frac{1}{2(\sigma^2)^2} \sum_{i=1}^{n}(x_i - \mu)^2 = 0.
\end{equation}

代入 $\hat{\mu} = \bar{x}$，解得：
\begin{equation}
    \boxed{\hat{\sigma}^2 = \frac{1}{n}\sum_{i=1}^{n}(x_i - \bar{x})^2}.
\end{equation}

\medskip
\textbf{第 4 步：验证极值点。}

可以验证二阶导数矩阵（Hessian）为负定，$\hat{\mu}, \hat{\sigma}^2$ 确为极大值点（略）。

\medskip
\textbf{注意：MLE 的 $\hat{\sigma}^2$ 分母是 $n$，不是 $n-1$！}

\begin{equation}
    \E[\hat{\sigma}^2_{\mathrm{MLE}}] = \frac{n-1}{n}\sigma^2 \neq \sigma^2.
\end{equation}

MLE 给出的方差估计是\emph{有偏的}（虽然渐近无偏）。
这就是为什么样本方差 $S^2$ 用 $n-1$ 做分母——修正了这个偏差。
\emph{MLE 追求最大化似然，不追求无偏性。}

\subsubsection{概念纠正：似然函数 $L(\theta)$ 究竟是谁的"概率"？}

以下对照你的理解逐条纠正。

\medskip
\textbf{你的第 2 点：}$f(x_i)$ 能知道每个观测值被取到的概率。

\emph{这正是最常见的误区。}
对\textbf{连续}分布，$f(x_i)$ 是\textbf{概率密度}，不是概率！
对于连续随机变量，单点概率恒为零：$P(X = x_i) = 0$。

但"密度"可以理解为"单位长度上的概率浓度"：$P(x_i < X < x_i + dx) \approx f(x_i) \cdot dx$。
虽然 $f(x_i)$ 本身不是概率（可以大于 1），但 $f(x_i)$ 越大意味着数据集中在 $x_i$ 附近的概率越高。

正因为单点概率为零，MLE 对连续分布用的是 PDF 的乘积而非概率的乘积——这是合理的，
因为要比较"哪组参数让观测数据更可能"，只需比较密度的相对大小。

\medskip
\textbf{你的第 3 点：}$L(\theta; x_i)$ 是在观测到 $\{x_i\}$ 的情况下参数为 $\theta$ 的概率。

\emph{这是错误的，而且是最关键的概念混淆。}

$L(\theta)$ \textbf{不等于} $P(\theta \mid \text{数据})$。两者的方向完全不同：

\begin{center}
\renewcommand{\arraystretch}{1.5}
\begin{tabular}{c|c|c}
\toprule
& \textbf{似然函数 $L(\theta)$} & \textbf{后验概率 $P(\theta \mid \text{数据})$} \\
\midrule
含义 & 参数为 $\theta$ 时，\emph{数据出现的可能性} & 观测到数据后，\emph{参数为 $\theta$ 的概率} \\
方向 & $\theta \to \text{数据的分布}$ & $\text{数据} \to \theta \text{ 的分布}$ \\
是否概率 & 否——$L(\theta)$ 不是 $\theta$ 上的概率分布 & 是——$\int P(\theta \mid \mathcal{D})\,d\theta = 1$ \\
所属框架 & \textbf{频率学派} & \textbf{贝叶斯学派} \\
\bottomrule
\end{tabular}
\end{center}

$L(\theta)$ 和 PDF 是同一个公式，只是视角不同——前者把数据固定、参数当变量，后者反之。
它不满足概率的归一化条件：$\int L(\theta)\,d\theta$ 一般不等于 1。

> \textbf{似然函数不是"参数的后验概率"，它是"固定数据后，参数与数据兼容程度的度量"。}

\medskip
\textbf{你的第 4 点：}为什么 $L$ 能和 PDF 挂钩？

\emph{答案：它们本质上是同一个数学表达式，只是"谁当变量"的视角切换。}

写出联合 PDF（$\theta$ 固定，$x$ 是变量）：
\begin{equation}
    f(x_1, \dots, x_n; \theta) = \prod_{i=1}^{n} f(x_i; \theta).
\end{equation}

现在做一次\textbf{视角翻转}——数据已经拿到了（$x_1,\dots,x_n$ 是固定常数），
问题是"哪组 $\theta$ 让这些数据最可能？"直接把上面的表达式看作 $\theta$ 的函数：
\begin{equation}
    L(\theta) = \prod_{i=1}^{n} f(x_i; \theta).
\end{equation}

\textbf{形式上二者一模一样，但数学身份不同：}
$f(x_1,\dots,x_n; \theta)$ 是 $x$ 上的概率密度（积分为 1）；
$L(\theta)$ 是 $\theta$ 上的函数（积分一般不为 1，不是密度）。

\begin{center}
\fbox{%
\begin{minipage}{0.88\textwidth}
\[
\underbrace{f(x_1,\dots,x_n; \theta)}_{\text{PDF：$\theta$ 固定，$x$ 变量}}
\quad\underbrace{\xrightarrow{\text{数据拿到手，$\theta$ 变参数}}}_{\text{同一个公式，视角翻转}}\quad
\underbrace{L(\theta; x_1,\dots,x_n)}_{\text{似然：$x$ 固定，$\theta$ 变量}}
\]
\end{minipage}}
\end{center}

\textbf{因此你的第 1 点方向正确，只需补充：}
我们有一组已观测到的数据 $x_1,\dots,x_n$。总体分布的形式已知（如正态），参数未知。目标是用观测数据反过来推断最合理的参数值——MLE 就是"哪组 $\theta$ 让手上这组数据出现的联合密度最大"。

\subsubsection{频率学派 vs. 贝叶斯学派：根本分歧}

MLE 和 MAP 分别代表统计推断的两大范式。理解它们的底层分歧，比记住公式更重要。

\medskip
\textbf{核心分歧：参数 $\theta$ 是什么？}

\begin{center}
\renewcommand{\arraystretch}{1.6}
\begin{tabular}{p{2.5cm}|p{5.5cm}|p{5.5cm}}
\toprule
& \textbf{频率学派 (Frequentist)} & \textbf{贝叶斯学派 (Bayesian)} \\
\midrule
$\theta$ 的本质
    & $\theta$ 是一个\emph{固定但未知的常数}。
      没有"$\theta$ 的分布"这个概念。
    & $\theta$ 是一个\emph{随机变量}，有自己的分布。
      在观测数据前，你对 $\theta$ 已有\emph{先验信念} $P(\theta)$。 \\
\midrule
概率的含义
    & 概率是\emph{长期频率}——无限次重复实验中事件发生的比例。
      "下次掷硬币正面的概率是 0.5" = 长期来看一半是正面。
    & 概率是\emph{信念程度}——你对一个命题的相信程度。
      "明天有 60\% 概率下雨" = 你对"明天下雨"有 60\% 的把握。 \\
\midrule
推断的目标
    & 给出 $\theta$ 的一个\emph{点估计}（MLE）或一个\emph{置信区间}。
    & 给出 $\theta$ 的\emph{整个后验分布} $P(\theta \mid \mathcal{D})$。 \\
\midrule
数据的角色
    & 数据是唯一的随机来源。参数固定，数据随机。 & 参数和数据都是随机变量。数据更新信念。 \\
\midrule
核心工具
    & MLE、假设检验、置信区间 & 贝叶斯公式、后验分布、可信区间 \\
\midrule
区间含义
    & \textbf{置信区间}：重复抽样 100 次，约 95 个区间包含 $\theta$。
      \emph{不能说}"这个区间有 95\% 概率包含 $\theta$"。
    & \textbf{可信区间}：给定数据，$\theta$ 落在此区间的\emph{概率}是 95\%。
      可以直接说"我们有 95\% 把握 $\theta$ 在区间内"。 \\
\bottomrule
\end{tabular}
\end{center}

\medskip
\textbf{用一个具体例子体会差异：}

你掷硬币 10 次，8 次正面。想知道硬币正面的概率 $p$。

\emph{频率学派（MLE）：}
\begin{itemize}
    \item $p$ 是一个固定的数，只是你不知道。
    \item $\hat{p} = 8/10 = 0.8$，这是 $p$ 的最佳点估计。
    \item "$p$ 的 95\% 置信区间是 $[0.49, 0.96]$" —— 如果重复这个实验无数次，95\% 的置信区间会包含真实的 $p$。
    \item \emph{你不能说}"$p$ 落在这个区间里的概率是 95\%"——$p$ 不是随机的。
\end{itemize}

\emph{贝叶斯学派（MAP / 后验）：}
\begin{itemize}
    \item $p$ 本身是不确定的——你相信 $p$ 可能接近 0.5（硬币通常公平），但也可能不是。
    \item 用 Beta(2,2) 作为先验（弱偏向公平），后验为 Beta(10,4)。
    \item 后验均值 $\frac{10}{14} \approx 0.714$（比 0.8 更保守，因为先验往回拉了一点）。
    \item "$p$ 的 95\% 可信区间是 $[0.46, 0.90]$" —— \emph{可以直接说}"有 95\% 的概率 $p$ 在这个区间里"。
\end{itemize}

\medskip
\textbf{MLE 和 MAP 在公式上的区别只有一项：}

\begin{equation}
    \underbrace{\hat{\theta}_{\mathrm{MLE}} = \argmax_\theta \log P(\mathcal{D} \mid \theta)}_{\text{只看数据}}
    \qquad
    \underbrace{\hat{\theta}_{\mathrm{MAP}} = \argmax_\theta \big[\log P(\mathcal{D} \mid \theta) + \log P(\theta)\big]}_{\text{数据 + 先验信念}}
\end{equation}

MAP 比 MLE 多了一项 $\log P(\theta)$ —— 这就是"你对参数已有的信念"。当 $P(\theta)$ 是均匀分布时，MAP = MLE。

\medskip
\textbf{现代 ML 中两类方法的分布：}

\begin{itemize}
    \item \textbf{频率派占主导的：} 大多数监督学习（线性回归、SVM、GBDT、深度学习训练）——目标就是找一个最好的 $\theta$，不关心它的分布。
    \item \textbf{贝叶斯派占主导的：} 不确定性量化（什么时候模型该说"我不知道"）、小样本学习、超参数调优（贝叶斯优化）、主题模型、变分推断。
    \item \textbf{融合趋势：} 现代 DL 中的 Dropout、BatchNorm 的推理可以看作变分推断；Ensemble 本质上是后验分布的蒙特卡洛近似。
\end{itemize}

\subsubsection{深层追问：把 $\theta$ 当常数 vs. 当变量，实际差别在哪？}

\emph{表面上}，MLE 和 MAP 都是求一个最优的 $\theta$，公式只差一项 $\log P(\theta)$。
但\textbf{底层的世界观分歧}远比公式大：

\medskip
\textbf{差别一：能否谈论"$\theta$ 的不确定性"。}

频率学派：你只能得到 $\hat{\theta}_{\mathrm{MLE}} = 0.8$，这是一个数。你可以构造置信区间，但\textbf{不能说} $\theta$ 落在这个区间的概率是 95\%——$\theta$ 不是随机的，它要么在里面，要么不在。95\% 是方法的长期成功率，不是本次推断的可信度。

贝叶斯学派：你得到整个后验分布 $P(\theta \mid \mathcal{D})$。你可以直接说"$\theta$ 有 95\% 的概率在 $[0.46, 0.90]$ 之间"，也可以量化"$\theta > 0.5$ 的概率是 97\%"。这在需要\textbf{决策}的场景中至关重要——如果有 10\% 的概率 $\theta$ 落在危险区域，你可能会采取完全不同的行动。

\medskip
\textbf{差别二：先验知识能不能进入推断。}

频率学派拒绝先验，坚持"让数据自己说话"。这在样本量大时没问题，但在\textbf{小样本}场景下会出荒谬结果：

\begin{itemize}
    \item 掷硬币 1 次，正面。MLE：$\hat{p} = 1.0$，认为硬币永远出正面。
    \item 掷硬币 3 次，3 次正面。MLE：$\hat{p} = 1.0$。但你的常识告诉你"硬币通常接近公平"。
\end{itemize}

贝叶斯学派用一个温和的先验（如 Beta(2,2)，相当于假设已经见过 2 正 2 反）把估计往回拉：
1 次正面 → 后验均值 $\frac{3}{5} = 0.6$（比 1.0 合理得多）。

这不是"主观偏见"，而是\textbf{用先验编码常识，防止小样本过拟合}。这在现代 ML 中对应正则化——Ridge 回归的 $L_2$ 惩罚等价于高斯先验，Lasso 的 $L_1$ 惩罚等价于拉普拉斯先验。

\medskip
\textbf{差别三：频率论的保证是"长期校准"而非"单次可信"。}

这是频率学派最强也最微妙的优势。

一个 95\% 置信区间的含义是：如果你把这个实验重复 100 次，每次算一个区间，大约 95 个会包含真实的 $\theta$。这是一种\textbf{过程保证}——方法本身在长期使用中是可靠的。

但贝叶斯 95\% 可信区间可能没有这个性质：在重复抽样中，贝叶斯区间可能只有 80\% 的次数覆盖真值（如果先验选得不好）。频率学派正是用这一点批评贝叶斯：你的区间说"95\% 概率包含 $\theta$"，但在客观世界中这个覆盖频率可能远低于 95\%。

\medskip
\textbf{为什么频率学派历史上排斥贝叶斯？三个核心原因：}

\begin{enumerate}[label=(\arabic*)]
    \item \textbf{先验的主观性（最根本的批评）}：\\
    "你的先验 $P(\theta)$ 从哪来？如果两个人用不同的先验，同一个数据得出不同的结论，科学还叫科学吗？"
    ——频率学派追求"客观"推断，只依赖数据，不依赖分析者的主观信念。

    \item \textbf{$\theta$ 经常确实不是随机的}：\\
    硬币的正面概率 $p$ 是一个物理常数（虽然你不知道具体值），把它当作随机变量建模在哲学上就是错的。
    频率学派的回答：我们不需要给 $p$ 一个分布——我们只需要一个估计方法和它的误差特性。

    \item \textbf{频率方法有客观的频率保证}：\\
    95\% 置信区间在长期中确实有 95\% 的覆盖率——这是数学定理，不依赖任何先验假设。
    贝叶斯可信区间没有这种频率主义校准——如果先验错了，区间也错了。
\end{enumerate}

\medskip
\textbf{现代视角：分歧在收敛，不在对立。}

今天的共识是：两种方法各有适用的场景，不是非此即彼的关系。

\begin{center}
\renewcommand{\arraystretch}{1.4}
\begin{tabular}{c|c|c}
\toprule
\textbf{场景} & \textbf{推荐方法} & \textbf{原因} \\
\midrule
$n$ 很大，只需点估计 & MLE（频率） & 简单、客观、有效。先验被数据淹没，两派结果接近 \\
\midrule
$n$ 很小，需要利用常识 & MAP / 后验（贝叶斯） & 纯靠数据会过拟合，先验起正则化作用 \\
\midrule
需要量化不确定性 & 后验分布（贝叶斯） & 频率区间不能说"概率"，决策需要真正的概率陈述 \\
\midrule
需要长期频率保证 & 频率方法 & 置信区间有客观的覆盖率保证 \\
\midrule
$p \gg n$（高维稀疏） & MAP（等价于正则化） & 频率派称为"加惩罚项"，贝叶斯派称为"加先验"，殊途同归 \\
\midrule
A/B 测试、临床试验 & 频率为主 & 监管机构要求频率主义的 $p$ 值和错误率控制 \\
\bottomrule
\end{tabular}
\end{center}

\textbf{一句话总结分歧的本质：}

\[
\text{频率派：}\;\theta \text{ 是客观的常数，推断 = 用随机数据去逼近固定真相。}
\]
\[
\text{贝叶斯派：}\;\theta \text{ 是主观的不确定性，推断 = 用数据更新你对真相的信念。}
\]

\subsection{最大后验估计 (MAP)}

引入先验分布 $P(\theta)$，最大化后验概率：
\begin{equation}
    \hat{\theta}_{\mathrm{MAP}} = \argmax_{\theta} P(\theta \mid \mathcal{D})
    = \argmax_{\theta} P(\mathcal{D} \mid \theta) P(\theta).
\end{equation}

\begin{remark}[第二个等号为什么成立？——贝叶斯公式 + 分母与 $\theta$ 无关]
贝叶斯公式为：
\begin{equation}
    P(\theta \mid \mathcal{D}) = \frac{P(\mathcal{D} \mid \theta) \, P(\theta)}{P(\mathcal{D})}.
\end{equation}

其中 $P(\mathcal{D}) = \int P(\mathcal{D} \mid \theta) P(\theta)\, d\theta$ 是数据的边缘概率。
\textbf{关键：$P(\mathcal{D})$ 与 $\theta$ 无关}——它是对所有可能的 $\theta$ 做了积分，积分完成后 $\theta$ 已消失。
在关于 $\theta$ 最大化时，$P(\mathcal{D})$ 是常数，可以直接扔掉：

\begin{equation}
    \argmax_{\theta} P(\theta \mid \mathcal{D})
    = \argmax_{\theta} \frac{P(\mathcal{D} \mid \theta) P(\theta)}{P(\mathcal{D})}
    = \argmax_{\theta} P(\mathcal{D} \mid \theta) P(\theta).
\end{equation}

\textbf{三步可视化：}

\[
\underbrace{P(\theta \mid \mathcal{D})}_{\text{后验}}
\;=\;
\frac{
    \overbrace{P(\mathcal{D} \mid \theta)}^{\text{似然}}
    \;\times\;
    \overbrace{P(\theta)}^{\text{先验}}
}{
    \underbrace{P(\mathcal{D})}_{\text{与 $\theta$ 无关，argmax 时可忽略}}
}
\;\xrightarrow{\argmax_\theta}\;
P(\mathcal{D} \mid \theta) \, P(\theta).
\]

\textbf{直觉：} 后验 $\propto$ 似然 $\times$ 先验。最大化后验 = 最大化（数据支持度 $\times$ 先验信念）。
先验 $P(\theta)$ 把 MLE 的纯数据驱动估计往你认为合理的区域"拉"了一把。
\end{remark}

当先验为均匀分布时，$P(\theta)$ 是常数，MAP 退化为 MLE。

% ============ §2 截尾样本的 MLE ============
\section{基于截尾样本的最大似然估计}

在可靠性试验和寿命试验中，由于时间和成本限制，试验往往在全部样品失效前终止，产生的数据称为\textbf{截尾样本}。

\subsection{定时截尾}

试验进行到规定时间 $t_0$ 时停止。假设 $n$ 个样品中有 $r$ 个在 $t_0$ 前失效，失效时间分别为 $t_1 \le t_2 \le \dots \le t_r \le t_0$。

似然函数：
\begin{equation}
    L(\theta) = \left[\prod_{i=1}^{r} f(t_i; \theta)\right] \cdot \big[1 - F(t_0; \theta)\big]^{n-r},
\end{equation}
其中 $f$ 为失效时间分布的 PDF，$F$ 为 CDF，$1-F(t_0)$ 表示一个样品在 $t_0$ 时仍未失效的概率（生存概率）。

\subsection{定数截尾}

试验进行到有 $r$ 个样品失效时停止，第 $r$ 个失效时间为 $t_r$。

似然函数：
\begin{equation}
    L(\theta) = \left[\prod_{i=1}^{r} f(t_i; \theta)\right] \cdot \big[1 - F(t_r; \theta)\big]^{n-r}.
\end{equation}

\subsection{指数分布截尾样本的 MLE}

设寿命 $T \sim \mathrm{Exp}(\lambda)$，即 $f(t) = \lambda e^{-\lambda t}$（$t \ge 0$），$F(t) = 1 - e^{-\lambda t}$。

定时截尾下的 MLE：
\begin{equation}
    \hat{\lambda} = \frac{r}{\sum_{i=1}^{r} t_i + (n-r)t_0}.
\end{equation}

% ============ §3 估计量的评选标准 ============
\section{估计量的评选标准}

\subsection{无偏性}

\begin{definition}[无偏估计量]
若估计量 $\hat{\theta}$ 满足 $\E[\hat{\theta}] = \theta$，则称 $\hat{\theta}$ 为 $\theta$ 的\textbf{无偏估计量}。
若 $\E[\hat{\theta}] \neq \theta$，其差 $\bias(\hat{\theta}) = \E[\hat{\theta}] - \theta$ 称为\textbf{偏差}。
\end{definition}

\begin{example}
$\bar{X}$ 是 $\mu$ 的无偏估计量；$S^2 = \frac{1}{n-1}\sum(X_i - \bar{X})^2$ 是 $\sigma^2$ 的无偏估计量，而 $\frac{1}{n}\sum(X_i - \bar{X})^2$ 是有偏的（渐近无偏）。
\end{example}

\subsection{有效性}

\begin{definition}[有效性]
设 $\hat{\theta}_1$ 与 $\hat{\theta}_2$ 都是 $\theta$ 的无偏估计量，若
\begin{equation}
    \Var(\hat{\theta}_1) \le \Var(\hat{\theta}_2),
\end{equation}
则称 $\hat{\theta}_1$ 比 $\hat{\theta}_2$ \textbf{有效}。

在所有无偏估计量中方差达到下界的称为\textbf{最小方差无偏估计量 (MVUE)}。
\end{definition}

\subsection{相合性（一致性）}

\begin{definition}[相合估计量]
若对任意 $\varepsilon > 0$，有
\begin{equation}
    \lim_{n \to \infty} P\big(|\hat{\theta}_n - \theta| < \varepsilon\big) = 1,
\end{equation}
即 $\hat{\theta}_n \xrightarrow{P} \theta$，则称 $\hat{\theta}_n$ 为 $\theta$ 的\textbf{相合估计量}（一致估计量）。
\end{definition}

\subsection{均方误差（综合评价）}

\begin{definition}[均方误差]
\begin{equation}
    \mse(\hat{\theta}) = \E\big[(\hat{\theta} - \theta)^2\big]
    = \Var(\hat{\theta}) + \big(\bias(\hat{\theta})\big)^2.
\end{equation}
MSE 同时考虑了方差和偏差，常用于有偏估计量的比较。
\end{definition}

% ============ §4 区间估计 ============
\section{区间估计}

\subsection{基本概念}

点估计给出了参数的一个具体值，但无法反映估计的精度。\textbf{区间估计}用一个随机区间来估计参数，并给出该区间包含参数真值的可信程度。

\begin{definition}[置信区间]
设 $\theta$ 为总体分布中的未知参数。对于给定的 $\alpha$（$0 < \alpha < 1$），若存在统计量
$\hat{\theta}_L = \hat{\theta}_L(X_1,\dots,X_n)$ 和 $\hat{\theta}_U = \hat{\theta}_U(X_1,\dots,X_n)$，使得
\begin{equation}
    P\big(\hat{\theta}_L < \theta < \hat{\theta}_U\big) = 1 - \alpha,
\end{equation}
则称随机区间 $(\hat{\theta}_L,\; \hat{\theta}_U)$ 为 $\theta$ 的\textbf{置信水平}为 $1 - \alpha$ 的\textbf{置信区间}。
\end{definition}

\begin{itemize}
    \item $1 - \alpha$：\textbf{置信水平}（置信度），常用 $0.90, 0.95, 0.99$
    \item $\hat{\theta}_L$：置信下限
    \item $\hat{\theta}_U$：置信上限
\end{itemize}

\textbf{精确解释：} 重复抽样 100 次，大约有 $100 \times (1-\alpha)$ 个区间包含 $\theta$。

\subsection{枢轴量法（构造置信区间的通用方法）}

\begin{enumerate}
    \item 找一个\textbf{枢轴量} $G = G(X_1,\dots,X_n; \theta)$，其分布完全已知且与 $\theta$ 无关；
    \item 对给定 $\alpha$，确定 $a, b$ 使 $P(a < G < b) = 1 - \alpha$；
    \item 解不等式 $a < G < b$ 得到 $\hat{\theta}_L < \theta < \hat{\theta}_U$。
\end{enumerate}

% ============ §5 正态总体均值与方差的区间估计 ============
\section{正态总体均值与方差的区间估计}

\subsection{单个正态总体均值的置信区间}

\subsubsection{$\sigma^2$ 已知时}

枢轴量：$\dfrac{\bar{X} - \mu}{\sigma / \sqrt{n}} \sim \mathcal{N}(0,1)$。

$\mu$ 的 $1-\alpha$ 置信区间：
\begin{equation}
    \boxed{\left(\bar{X} - z_{\alpha/2} \cdot \frac{\sigma}{\sqrt{n}},\;
                 \bar{X} + z_{\alpha/2} \cdot \frac{\sigma}{\sqrt{n}}\right)}.
\end{equation}

\textbf{常用分位点：} $z_{0.025} = 1.96$（$\alpha=0.05$），$z_{0.005} = 2.576$（$\alpha=0.01$）。

\subsubsection{$\sigma^2$ 未知时}

枢轴量：$\dfrac{\bar{X} - \mu}{S / \sqrt{n}} \sim t(n-1)$。

$\mu$ 的 $1-\alpha$ 置信区间：
\begin{equation}
    \boxed{\left(\bar{X} - t_{\alpha/2}(n-1) \cdot \frac{S}{\sqrt{n}},\;
                 \bar{X} + t_{\alpha/2}(n-1) \cdot \frac{S}{\sqrt{n}}\right)}.
\end{equation}

\subsection{单个正态总体方差的置信区间}

枢轴量：$\dfrac{(n-1)S^2}{\sigma^2} \sim \chi^2(n-1)$。

$\sigma^2$ 的 $1-\alpha$ 置信区间：
\begin{equation}
    \boxed{\left(\frac{(n-1)S^2}{\chi^2_{\alpha/2}(n-1)},\;
                 \frac{(n-1)S^2}{\chi^2_{1-\alpha/2}(n-1)}\right)}.
\end{equation}

注意：$\chi^2$ 分布不对称，需双侧分位点。

\subsection{两个正态总体均值差的置信区间}

\subsubsection{$\sigma_1^2, \sigma_2^2$ 已知}

\begin{equation}
    \boxed{\left((\bar{X} - \bar{Y}) - z_{\alpha/2} \sqrt{\frac{\sigma_1^2}{n_1} + \frac{\sigma_2^2}{n_2}},\;
                 (\bar{X} - \bar{Y}) + z_{\alpha/2} \sqrt{\frac{\sigma_1^2}{n_1} + \frac{\sigma_2^2}{n_2}}\right)}.
\end{equation}

\subsubsection{$\sigma_1^2 = \sigma_2^2 = \sigma^2$ 但 $\sigma^2$ 未知}

枢轴量涉及 $S_w$（合并方差）和 $t(n_1+n_2-2)$：
\begin{equation}
    \boxed{\left((\bar{X} - \bar{Y}) \pm t_{\alpha/2}(n_1+n_2-2) \cdot S_w \sqrt{\frac{1}{n_1} + \frac{1}{n_2}}\right)}.
\end{equation}

\subsection{两个正态总体方差比的置信区间}

枢轴量：$\dfrac{S_1^2 / \sigma_1^2}{S_2^2 / \sigma_2^2} \sim F(n_1-1, n_2-1)$。

$\dfrac{\sigma_1^2}{\sigma_2^2}$ 的 $1-\alpha$ 置信区间：
\begin{equation}
    \boxed{\left(\frac{S_1^2}{S_2^2} \cdot \frac{1}{F_{\alpha/2}(n_1-1,\, n_2-1)},\;
                 \frac{S_1^2}{S_2^2} \cdot \frac{1}{F_{1-\alpha/2}(n_1-1,\, n_2-1)}\right)}.
\end{equation}

% ============ §6 (0-1)分布参数的区间估计 ============
\section{(0-1) 分布参数的区间估计}

设总体 $X \sim B(1, p)$，即 $P(X=1) = p$，$P(X=0) = 1-p$。

由中心极限定理，当 $n$ 充分大时：
\begin{equation}
    \frac{\bar{X} - p}{\sqrt{p(1-p)/n}} \xrightarrow{d} \mathcal{N}(0,1).
\end{equation}

$p$ 的近似 $1-\alpha$ 置信区间（解二次不等式）：
\begin{equation}
    \boxed{\frac{1}{1 + \frac{z_{\alpha/2}^2}{n}}
           \left(\hat{p} + \frac{z_{\alpha/2}^2}{2n}
           \pm z_{\alpha/2} \sqrt{\frac{\hat{p}(1-\hat{p})}{n} + \frac{z_{\alpha/2}^2}{4n^2}}\right)},
\end{equation}
其中 $\hat{p} = \bar{X} = \frac{1}{n}\sum X_i$。

当 $n$ 很大且 $\hat{p}$ 不接近 0 或 1 时，可用简化公式：
\begin{equation}
    \boxed{\left(\hat{p} - z_{\alpha/2} \sqrt{\frac{\hat{p}(1-\hat{p})}{n}},\;
                 \hat{p} + z_{\alpha/2} \sqrt{\frac{\hat{p}(1-\hat{p})}{n}}\right)}.
\end{equation}

% ============ §7 单侧置信区间 ============
\section{单侧置信区间}

在某些实际问题中，我们只关心参数的上限（不超过多少）或下限（不低于多少），此时使用\textbf{单侧置信区间}。

\begin{definition}[单侧置信区间]
\begin{itemize}
    \item \textbf{单侧置信上限}：若 $P(\theta < \hat{\theta}_U) = 1 - \alpha$，则 $(-\infty,\; \hat{\theta}_U)$ 为单侧置信区间，$\hat{\theta}_U$ 为单侧置信上限；
    \item \textbf{单侧置信下限}：若 $P(\theta > \hat{\theta}_L) = 1 - \alpha$，则 $(\hat{\theta}_L,\; +\infty)$ 为单侧置信区间，$\hat{\theta}_L$ 为单侧置信下限。
\end{itemize}
\end{definition}

\textbf{构造方法：} 与双侧类似，但将 $\alpha/2$ 换为 $\alpha$：

\begin{example}[正态总体均值的单侧置信区间（$\sigma^2$ 未知）]
\begin{itemize}
    \item 单侧上限：$\mu < \bar{X} + t_{\alpha}(n-1) \cdot \dfrac{S}{\sqrt{n}}$
    \item 单侧下限：$\mu > \bar{X} - t_{\alpha}(n-1) \cdot \dfrac{S}{\sqrt{n}}$
\end{itemize}
\end{example}

\begin{example}[正态总体方差的单侧置信区间]
$\sigma^2$ 的单侧上限：
\begin{equation}
    \sigma^2 < \frac{(n-1)S^2}{\chi^2_{1-\alpha}(n-1)}.
\end{equation}
如材料的强度只需关注下限，寿命的方差只需关注上限。
\end{example}

% ============================================================
%  第 N 章：假设检验
% ============================================================
